% Blueprint: Zero-Sum Problems, Extremal Set Theory, and P ≠ NP
%
% This document serves as the mathematical blueprint for the Lean 4
% formalization. Each definition and theorem is linked to its Lean
% declaration via \lean{} and marked with \leanok when fully formalized.
% Dependencies are tracked via \uses{} for the dependency graph.

% Common macros for both web and print output
\newcommand{\ZZ}{\mathbb{Z}}
\newcommand{\NN}{\mathbb{N}}
\newcommand{\PP}{\mathsf{P}}
\newcommand{\NP}{\mathsf{NP}}
\DeclareMathOperator{\ZMod}{ZMod}
\newcommand{\Dav}{\mathrm{D}}
\newcommand{\EGZ}{\mathrm{EGZ}}


\title{Zero-Sum Problems, Extremal Set Theory, and $\PP \neq \NP$}
\author{zpnenp project}

\chapter{Foundations}

\section{Subset Sum}

\begin{definition}[Subset Sum]
  \label{def:subset_sum}
  \lean{SubsetSum}
  \leanok
  Let $S$ be a finite set of integers and $t \in \ZZ$. The \emph{Subset Sum}
  predicate $\mathrm{SubsetSum}(S, t)$ holds if there exists $S' \subseteq S$
  such that $\sum_{x \in S'} x = t$.
\end{definition}

\begin{definition}[Zero-Sum Subset]
  \label{def:subset_sum_zero}
  \lean{SubsetSumZero}
  \leanok
  \uses{def:subset_sum}
  $\mathrm{SubsetSumZero}(S)$ holds if there exists a nonempty $S' \subseteq S$
  with $\sum_{x \in S'} x = 0$.
\end{definition}

\begin{definition}[Achievable Subset Sums]
  \label{def:subset_sums}
  \lean{subsetSums}
  \leanok
  \uses{def:subset_sum}
  The set of all achievable subset sums of $S$:
  \[
    \mathrm{subsetSums}(S) = \left\{ \sum_{x \in S'} x \;\middle|\; S' \subseteq S \right\}.
  \]
\end{definition}

\begin{theorem}[Subset Sum at Zero]
  \label{thm:subset_sum_zero}
  \lean{subsetSum_zero}
  \leanok
  \uses{def:subset_sum}
  For any finite set $S$ of integers, $\mathrm{SubsetSum}(S, 0)$ holds
  (witnessed by the empty subset).
\end{theorem}

\begin{theorem}[Monotonicity]
  \label{thm:subset_sum_mono}
  \lean{subsetSum_mono}
  \leanok
  \uses{def:subset_sum}
  If $S \subseteq T$ and $\mathrm{SubsetSum}(S, t)$, then $\mathrm{SubsetSum}(T, t)$.
\end{theorem}

\begin{theorem}[Singleton]
  \label{thm:subset_sum_singleton}
  \lean{subsetSum_singleton}
  \leanok
  \uses{def:subset_sum}
  $\mathrm{SubsetSum}(\{a\}, t)$ if and only if $t = 0$ or $t = a$.
\end{theorem}

\begin{theorem}[Counting Bound]
  \label{thm:card_subset_sums_le}
  \lean{card_subsetSums_le}
  \leanok
  \uses{def:subset_sums}
  $|\mathrm{subsetSums}(S)| \leq 2^{|S|}$.
\end{theorem}

\section{Modular Subset Sum and EGZ}

\begin{definition}[Modular Zero-Sum]
  \label{def:mod_subset_sum_zero}
  \lean{ModSubsetSumZero}
  \leanok
  $\mathrm{ModSubsetSumZero}(S, n)$ holds if there exists a nonempty
  $S' \subseteq S$ with $n \mid \sum_{x \in S'} x$.
\end{definition}

\begin{theorem}[EGZ for Finsets]
  \label{thm:egz_finset}
  \lean{egz_finset}
  \leanok
  \uses{def:mod_subset_sum_zero}
  If $|S| \geq 2n - 1$ and $n > 0$, then there exists $T \subseteq S$
  with $|T| = n$ and $n \mid \sum_{x \in T} x$.
\end{theorem}

\begin{theorem}[EGZ Implies Modular Zero-Sum]
  \label{thm:egz_mod}
  \lean{egz_implies_modSubsetSumZero}
  \leanok
  \uses{thm:egz_finset, def:mod_subset_sum_zero}
  If $|S| \geq 2n - 1$ and $n > 1$, then $\mathrm{ModSubsetSumZero}(S, n)$.

  \textbf{Significance:} The adversary cannot construct a set of $\geq 2n-1$
  integers that avoids having an $n$-element subset with sum divisible by~$n$.
  Existence is forced by combinatorial structure.
\end{theorem}

\section{The Davenport Constant}

\begin{definition}[Zero-Sum Free]
  \label{def:zero_sum_free}
  \lean{ZeroSumFree}
  \leanok
  A multiset $s$ over an additive group is \emph{zero-sum free} if no nonempty
  submultiset of $s$ sums to zero.
\end{definition}

\begin{theorem}[Sum of Replicate]
  \label{thm:sum_replicate}
  \lean{multiset_sum_replicate}
  \leanok
  The sum of $k$ copies of $a$ equals $k \cdot a$ (i.e., $k \smul a$).
\end{theorem}

\begin{theorem}[Lower Bound: $\Dav(\ZZ/n\ZZ) \geq n$]
  \label{thm:davenport_lower}
  \lean{zeroSumFree_replicate_one}
  \leanok
  \uses{def:zero_sum_free, thm:sum_replicate}
  For $n > 1$, the multiset of $(n-1)$ copies of $1$ in $\ZZ/n\ZZ$ is
  zero-sum free. Therefore $\Dav(\ZZ/n\ZZ) \geq n$.

  \begin{proof}
    Any nonempty submultiset consists of $k$ copies of $1$ for $1 \leq k \leq n-1$.
    Its sum is $k \in \ZZ/n\ZZ$. Since $1 \leq k \leq n-1$, we have $n \nmid k$,
    so $k \neq 0$ in $\ZZ/n\ZZ$.
  \end{proof}

  \textbf{Key insight:} The extremal instance (longest zero-sum free multiset)
  is \emph{maximally structured}---it consists entirely of copies of a single
  element. Hard instances are not random; they are rigid.
\end{theorem}

\begin{theorem}[Upper Bound: $\Dav(\ZZ/n\ZZ) \leq n$]
  \label{thm:davenport_upper}
  \lean{davenport_upper}
  \uses{def:zero_sum_free}
  For $n > 0$, every multiset of $n$ elements in $\ZZ/n\ZZ$ contains a
  nonempty submultiset summing to zero.

  \begin{proof}[Proof sketch]
    Let $a_1, \ldots, a_n \in \ZZ/n\ZZ$. Consider the prefix sums
    $s_0 = 0, s_1 = a_1, s_2 = a_1 + a_2, \ldots, s_n = a_1 + \cdots + a_n$.
    These are $n+1$ elements of $\ZZ/n\ZZ$, which has $n$ elements. By the
    pigeonhole principle, $s_i = s_j$ for some $i < j$. Then
    $a_{i+1} + \cdots + a_j = s_j - s_i = 0$.
  \end{proof}
\end{theorem}

\begin{theorem}[$\Dav(\ZZ/n\ZZ) = n$]
  \label{thm:davenport_zmod}
  \lean{davenport_ZMod}
  \uses{thm:davenport_lower, thm:davenport_upper}
  For $n > 1$:
  \begin{enumerate}
    \item Every multiset of $n$ elements in $\ZZ/n\ZZ$ has a nonempty
          zero-sum submultiset.
    \item There exists a multiset of $n-1$ elements in $\ZZ/n\ZZ$ with
          no nonempty zero-sum submultiset.
  \end{enumerate}
\end{theorem}

\chapter{Structural Theory}

\section{Inverse Zero-Sum Theorems}

\begin{theorem}[Prefix Sums Injective]
  \label{thm:prefix_sums_injective}
  \lean{ZeroSumFree.prefix_sums_injective}
  \leanok
  \uses{def:zero_sum_free}
  If a multiset (as a list) is zero-sum free, its prefix sums are all distinct.
\end{theorem}

\begin{theorem}[All Elements Equal]
  \label{thm:all_eq}
  \lean{ZeroSumFree.all_eq}
  \leanok
  \uses{def:zero_sum_free, thm:prefix_sums_injective}
  In a zero-sum free multiset of size $n-1$ in $\ZZ/n\ZZ$, all elements
  are equal. (Proved via adjacent-swap prefix sum argument.)
\end{theorem}

\begin{theorem}[Inverse Davenport for $\ZZ/n\ZZ$]
  \label{thm:inverse_davenport}
  \lean{inverse_davenport}
  \leanok
  \uses{def:zero_sum_free, thm:all_eq}
  The zero-sum free multisets of maximal length $n-1$ in $\ZZ/n\ZZ$ are
  precisely those consisting of $n-1$ copies of a single element $g$
  with $\gcd(g, n) = 1$.
\end{theorem}

\begin{definition}[EGZ-Free]
  \label{def:egz_free}
  \lean{EGZFree}
  \leanok
  A multiset is \emph{EGZ-free} if no submultiset of size $n$ sums to zero.
\end{definition}

\begin{theorem}[Extremal EGZ-Free (Forward)]
  \label{thm:egz_extremal_free}
  \lean{egzExtremal_free}
  \leanok
  \uses{def:egz_free}
  The multiset of $(n-1)$ copies of $a$ and $(n-1)$ copies of $b$,
  with $a - b$ a unit, is EGZ-free.
\end{theorem}

\begin{theorem}[At Most Two Values]
  \label{thm:at_most_two_values}
  \lean{EGZFree.at_most_two_values}
  \uses{def:egz_free, thm:inverse_davenport}
  An EGZ-free multiset of size $2n-2$ has at most $2$ distinct values.
  (Sorry: Gao's theorem subcase for $n \geq 5$.)
\end{theorem}

\begin{theorem}[Inverse EGZ]
  \label{thm:inverse_egz}
  \lean{inverse_egz}
  \uses{thm:egz_extremal_free, thm:at_most_two_values}
  The multisets of length $2n-2$ in $\ZZ/n\ZZ$ that do not contain an
  $n$-element zero-sum submultiset are precisely those consisting of
  $(n-1)$ copies of element $a$ and $(n-1)$ copies of element $b$,
  with $a - b$ a unit in $\ZZ/n\ZZ$.
\end{theorem}

\section{Sumset Bounds}

\begin{theorem}[Cauchy--Davenport]
  \label{thm:cauchy_davenport}
  \leanok
  For $A, B \subseteq \ZZ/p\ZZ$ with $p$ prime:
  $|A + B| \geq \min(p,\; |A| + |B| - 1)$.
  (From Mathlib: \texttt{ZMod.cauchy\_davenport}.)
\end{theorem}

\begin{theorem}[Doubling in $\ZZ/p\ZZ$]
  \label{thm:doubling_gt_one}
  \lean{addDoubling_gt_one_of_small}
  \leanok
  \uses{thm:cauchy_davenport}
  For $A \subseteq \ZZ/p\ZZ$ with $2 \leq |A| \leq p-1$:
  $|A| < |A + A|$.
\end{theorem}

\begin{theorem}[Iterated Sumset Growth]
  \label{thm:iterated_sumset}
  \lean{iterated_sumset_growth}
  \leanok
  \uses{thm:cauchy_davenport}
  $|kA| \geq \min(p,\; k|A| - k + 1)$.
\end{theorem}

\begin{theorem}[Freiman for $\ZZ/p\ZZ$]
  \label{thm:freiman}
  \lean{freiman_ZMod}
  \uses{thm:cauchy_davenport}
  If $|A + A| \leq K|A|$ and $|A| \leq p/K$, then $A$ is contained in
  an arithmetic progression of length $\leq K^2 |A|$.
  (Sorry: deep theorem.)
\end{theorem}

\section{Sum-Product Phenomenon}

\begin{theorem}[Sum-Product Growth]
  \label{thm:sum_product}
  \lean{bourgain_katz_tao}
  \leanok
  \uses{thm:doubling_gt_one}
  For $A \subseteq \ZZ/p\ZZ$ with $2 \leq |A| \leq p/2$:
  $|A| < \max(|A + A|, |A \cdot A|)$.
\end{theorem}

\section{Density Phase Transitions}

\begin{definition}[Density Regimes]
  \label{def:density}
  \lean{SubsetSumInstance.isHighDensity}
  \leanok
  A Subset Sum instance has \emph{high density} if $2^n > n \cdot M + 1$,
  \emph{low density} if $M > 2^{n^2}$, and \emph{critical density} otherwise.
\end{definition}

\begin{theorem}[Pigeonhole Collision]
  \label{thm:pigeonhole_collision}
  \lean{pigeonhole_collision}
  \leanok
  \uses{def:density}
  At high density, two distinct subsets have the same sum.
\end{theorem}

\chapter{Complexity Bridge}

\section{Subset Sum in NP}

\begin{theorem}[Subset Sum $\in \NP$]
  \label{thm:subset_sum_np}
  \lean{subsetSum_in_NP}
  \leanok
  \uses{def:subset_sum}
  $\mathrm{SubsetSum}(S, t) \iff \exists S' \in \mathcal{P}(S),\; \sum_{x \in S'} x = t$.
  Verification is decidable (polynomial-time checkable).
\end{theorem}

\begin{theorem}[Subset Sum Decidable]
  \label{thm:subset_sum_decidable}
  \lean{subsetSum_decidable}
  \leanok
  \uses{thm:subset_sum_np}
  The Subset Sum predicate has a \texttt{Decidable} instance.
\end{theorem}

\section{Query Complexity}

\begin{theorem}[$\Omega(n)$ Query Lower Bound]
  \label{thm:query_lower_bound}
  \lean{query_lower_bound_finset}
  \leanok
  \uses{def:subset_sum}
  The adversary can produce two Finsets differing by one element
  with different Subset Sum answers: $\{0, 1\}$ vs $\{0\}$ with target~$1$.
  Any algorithm must read all elements.
\end{theorem}

\section{The Adversary Framework}

\begin{theorem}[Adversary NO Instances]
  \label{thm:adversary_no}
  \lean{adversary_no_instances}
  \leanok
  \uses{thm:inverse_davenport}
  At the Davenport threshold (size $n-1$), the adversary's NO instances
  are exactly the replicate-of-unit instances.
\end{theorem}

\section{Structure vs.\ Computation}

\begin{theorem}[Structure--Computation Dichotomy]
  \label{thm:struct_comp_dichotomy}
  \lean{structure_vs_computation_dichotomy}
  \leanok
  \uses{thm:doubling_gt_one}
  For non-trivial $A \subseteq \ZZ/p\ZZ$: either $|A+A| > |A|$ (growing sumset)
  or $|A+A| \leq K|A|$ (Freiman structure). The adversary cannot avoid both.
\end{theorem}

\begin{theorem}[Adversary Full Dichotomy]
  \label{thm:adversary_full}
  \lean{adversary_full_dichotomy}
  \leanok
  \uses{thm:struct_comp_dichotomy, thm:sum_product}
  Every non-trivial instance simultaneously has growing sumsets AND large
  sum-product quantity: $|A| < |A+A|$ and $|A| < \max(|A+A|, |A \cdot A|)$.
\end{theorem}

\section{Barrier Analysis}

\begin{theorem}[Structural Theory is Combinatorial]
  \label{thm:barriers}
  \lean{structural_theory_is_combinatorial}
  \leanok
  \uses{thm:inverse_davenport}
  The structural theorems are purely algebraic (no computational content).

  \textbf{Barrier assessment:}
  \begin{itemize}
    \item Relativization: pre-computational (N/A)
    \item Natural proofs: potentially avoided (structural properties are rare)
    \item Algebrization: potentially hit (theory is algebraic)
  \end{itemize}
\end{theorem}

\chapter{Open Problems}

\begin{theorem}[$\PP \neq \NP$]
  \label{thm:p_ne_np}
  \uses{thm:inverse_davenport, thm:inverse_egz, thm:struct_comp_dichotomy, thm:subset_sum_np, thm:barriers}
  \textbf{Goal.} No polynomial-time algorithm solves Subset Sum in the
  worst case.

  \textbf{Strategy:} Use the structural characterization of hard instances
  from zero-sum theory to show that any polynomial-time algorithm must fail
  on instances at critical density with intermediate structure.

  \textbf{Status:} The structural theory constrains WHAT hard instances
  look like. The gap is showing WHY this structure implies computational
  hardness. The barrier analysis identifies algebrization as the main obstacle.
\end{theorem}
